% Mục nội dung thuộc chương 2 của bài giảng gốc; không phải chương độc lập.

\section{Mục tiêu học tập và vị trí trong kiến trúc IoT}
Hoàn thành chương này, người học cần: định nghĩa được cảm biến; mô tả đường đi của đại lượng đo đến dữ liệu số; phân biệt cảm biến tích cực/thụ động và đầu ra tương tự/số; nhận diện các cảm biến thường gặp; tính toán giá trị ADC và lựa chọn giao tiếp phù hợp; thiết kế sơ bộ nút thu thập dữ liệu với chú ý về nguồn, sai số và hiệu chuẩn.

Theo slide bài giảng, kiến trúc IoT được hình dung qua ba tầng gần gũi: \textbf{Things / IoT Devices} thu thập dữ liệu vật lý; \textbf{Edge} xử lý gần nơi phát sinh dữ liệu; và \textbf{Cloud} lưu trữ, phân tích, quản lý thông qua kết nối mạng. Phần cứng nút cảm biến là cầu nối giữa thế giới vật lý và luồng thông tin số.
\sourcepages{3--4, 16--17}

\begin{figure}[htbp]\centering
\begin{tikzpicture}[node distance=7mm, font=\small\sffamily]
\node[block] (world) {Môi trường\\$m$: nhiệt, áp, sáng, khí};
\node[block,right=of world] (trans) {Phần tử nhạy\\Biến đổi vật lý};
\node[block,right=of trans] (sig) {Điều hòa tín hiệu\\Khuếch đại / lọc};
\node[block,below=of sig] (adc) {ADC / giao tiếp số\\Số hóa mẫu đo};
\node[block,left=of adc] (mcu) {MCU / Edge\\Đọc, xử lý, hiệu chuẩn};
\draw[arrow] (world)--(trans); \draw[arrow](trans)--(sig); \draw[arrow](sig)--(adc); \draw[arrow](adc)--(mcu);
\draw[arrow] (mcu.west)--++(-9mm,0) node[left,font=\scriptsize,align=center] {Gửi mạng\\hoặc điều khiển};
\end{tikzpicture}
\caption{Chuỗi thu thập dữ liệu cảm biến, sơ đồ biên soạn mới.}
\end{figure}

\section{Khái niệm và đặc tính chuyển đổi}
\hwterm{Cảm biến}{sensor} là phần tử hoặc thiết bị nhận một đại lượng vật lý, hóa học, sinh học và tạo ra tín hiệu đo có thể đọc, xử lý hoặc truyền đi. Đại lượng được đo ký hiệu $m$ (\emph{measurand}); tín hiệu đầu ra là $s$; quan hệ chuyển đổi được viết
\begin{equation}\boxed{s=f(m).}\label{eq:transfer}\end{equation}
Trường hợp gần tuyến tính trong một miền làm việc: $s=km+b$. Hệ số $k=\Delta s/\Delta m$ được gọi là \emph{độ nhạy} nếu các điều kiện đo phù hợp. Tín hiệu có thể là điện áp, dòng điện, điện trở, điện dung, tần số, xung hoặc chuỗi bit qua bus số. Trên thực tế, nhiều cảm biến cần \emph{mạch kích thích} và mạch điều hòa để tạo ra đầu ra phù hợp với bộ xử lý.
\sourcepages{4}

\subsection{Các đại lượng ảnh hưởng chất lượng phép đo (mở rộng kỹ thuật)}
\begin{description}[leftmargin=0mm,labelindent=0mm,style=nextline]
\item[Khoảng đo (range)] Khoảng đại lượng mà thiết bị cho phép đo theo công bố của nhà sản xuất; vượt khoảng đo có thể bão hòa hoặc hư hỏng.
\item[Độ phân giải (resolution)] Thay đổi nhỏ nhất có thể phân biệt tại đầu ra; độ phân giải số không đồng nghĩa với độ chính xác.
\item[Độ chính xác (accuracy)] Mức độ gần giá trị đo với giá trị chuẩn; thường được nêu theo $\pm$ đơn vị hoặc phần trăm.
\item[Độ lặp lại (repeatability) và độ chính xác lặp (precision)] Mức phân tán kết quả trong điều kiện đo giống nhau; kết quả lặp tốt vẫn có thể lệch chuẩn.
\item[Độ nhạy, độ tuyến tính, độ trễ] Mô tả mức đáp ứng, sai lệch so với mô hình tuyến tính và sự phụ thuộc vào chiều tăng/giảm đại lượng.
\item[Thời gian đáp ứng, dải thông] Giới hạn tốc độ phản ứng, ảnh hưởng trực tiếp đến tần số lấy mẫu phù hợp.
\item[Độ trôi (drift)] Biến đổi điểm không hoặc độ nhạy do nhiệt độ, lão hóa và môi trường. Phải xem xét hiệu chuẩn định kỳ.
\end{description}

\section{Phân loại cảm biến}
Slide cung cấp ba hướng phân loại: theo \emph{nguồn năng lượng/tác động kích thích}, theo \emph{đại lượng hay hiện tượng chuyển đổi} và theo \emph{dạng đầu ra}. Không nên hiểu đây là ba cách phân loại loại trừ nhau: một cảm biến nhiệt điện trở có thể đồng thời được gọi là cảm biến nhiệt, cảm biến điện trở và cảm biến đầu ra analog khi ghép với mạch đọc.
\sourcepages{7--12}

\subsection{Cảm biến tích cực và thụ động}
Bài giảng mô tả cảm biến \hwterm{tích cực}{active} là cảm biến cần năng lượng kích thích bên ngoài, ví dụ radar/sonar chủ động phát sóng và đo phản xạ; cảm biến \hwterm{thụ động}{passive} là cảm biến không cần chủ động phát tín hiệu đo, ví dụ cặp nhiệt điện tự phát sinh sức điện động từ chênh lệch nhiệt độ.

\begin{notebox}
\textbf{Phân biệt quan trọng:} Trong tài liệu chuyên ngành, ``active/passive'' không phải lúc nào cũng được định nghĩa thống nhất. Việc cảm biến \emph{không chủ động phát sóng} khác với việc \emph{không cần nguồn cấp}. Camera thụ động thu ánh sáng môi trường nhưng vẫn cần điện để hoạt động. Một cảm biến áp suất kiểu điện trở/điện dung thường cần mạch kích thích, trong khi phần tử áp điện có thể tự tạo điện tích. Khi trả lời câu hỏi, nêu rõ tiêu chí phân loại được sử dụng.
\end{notebox}

\begin{table}[htbp]\centering\small
\caption{Đối chiếu nguyên lý theo cách mô tả của bài giảng}
\begin{tabularx}{\textwidth}{p{31mm}XX}
\toprule
Tiêu chí & Tích cực (phát tín hiệu dò) & Thụ động (khai thác tín hiệu sẵn có)\\\midrule
Nguyên lý & Phát siêu âm, vô tuyến hoặc ánh sáng rồi đo đáp ứng & Thu tín hiệu từ đại lượng tự nhiên hoặc chuyển đổi trực tiếp\\
Ví dụ & Radar đo khoảng cách; siêu âm HC-SR04 & Cặp nhiệt điện; quang trở nhận ánh sáng\\
Ưu điểm & Chủ động phép đo, ít phụ thuộc một số nguồn kích thích bên ngoài & Có thể đơn giản, ít ảnh hưởng tới đối tượng\\
Lưu ý & Công suất phát, phản xạ và nhiễu đồng kênh & Không đồng nghĩa toàn bộ thiết bị không dùng điện\\\bottomrule
\end{tabularx}
\end{table}

\subsection{Tín hiệu tương tự và tín hiệu số}
\hwterm{Tín hiệu tương tự}{analog} thay đổi liên tục theo thời gian/giá trị (trong mô hình lý tưởng). \hwterm{Tín hiệu số}{digital} biểu diễn thông tin bằng các mức rời rạc; tín hiệu logic vật lý thường biểu diễn mức LOW/HIGH (0/1), còn dữ liệu đo số có thể mã hóa hàng chục bit và nhiều giá trị, không bị giới hạn ở hai giá trị của đại lượng đo.

\begin{figure}[htbp]\centering
\begin{tikzpicture}[>=Latex,scale=.92,font=\small\sffamily]
\draw[->] (0,0)--(5.4,0) node[right]{Thời gian}; \draw[->] (0,-.2)--(0,2.2) node[above]{V};
\draw[very thick,tealiot,domain=0:5,samples=100] plot (\x,{1+.75*sin(180*\x/1.25)});
\node at (2.6,-.55){Tương tự: biến thiên liên tục};
\begin{scope}[xshift=7cm]
\draw[->] (0,0)--(5.4,0) node[right]{Thời gian}; \draw[->] (0,-.2)--(0,2.2) node[above]{Mức logic};
\draw[ultra thick,orangeiot] (0,0)--(.7,0)--(.7,1.6)--(1.7,1.6)--(1.7,0)--(2.5,0)--(2.5,1.6)--(3.8,1.6)--(3.8,0)--(5,0);
\node at (2.6,-.55){Số: mức thấp / cao};
\end{scope}
\end{tikzpicture}
\caption{Hai dạng tín hiệu thường gặp ở đầu vào vi điều khiển.}
\end{figure}

Một \emph{cảm biến analog} có thể xuất điện áp tỷ lệ đại lượng đo (hoặc thay đổi điện trở cần mạch phân áp). Bộ điều khiển cần ADC để chuyển thành mã số. Một \emph{cảm biến số} tích hợp ADC, xử lý/hiệu chuẩn hoặc bộ giao tiếp số, truyền kết quả qua I\textsuperscript{2}C, SPI, UART, 1-Wire; cũng có thiết bị chỉ tạo xung logic. Do đó nhận định trong slide rằng cảm biến số ``có dải đo thấp'' không phải tính chất phổ quát: phải đối chiếu thông số của \emph{từng model}.
\sourcepages{9--12}

\subsection{Phân loại theo đại lượng và cơ chế}
\begin{table}[htbp]\centering\footnotesize
\caption{Các cảm biến minh họa trong bài giảng và công dụng}
\begin{tabularx}{\textwidth}{p{32mm}p{40mm}X}
\toprule
Nhóm & Ví dụ trong slide & Đại lượng / ứng dụng\\\midrule
Nhiệt độ & Temperature sensor; cặp nhiệt điện & Giám sát nhiệt độ, bảo vệ thiết bị\\
Ẩm & Humidity, soil moisture & Độ ẩm không khí hoặc độ ẩm đất; không dùng lẫn hai đơn vị đo\\
Quang & LDR, photodiode, IR optical & Cường độ ánh sáng, phát hiện vật/đường line\\
Khoảng cách & Ultrasonic, IR proximity & Khoảng cách, nhận diện chướng ngại\\
Chuyển động & Accelerometer, gyroscope, PIR & Gia tốc, tốc độ góc, phát hiện chuyển động hồng ngoại\\
Môi trường & Gas, CO$_2$, rain, alcohol & Chỉ thị khí hoặc lượng mưa; phân biệt cảm biến khí tổng hợp và CO$_2$ chuyên dụng\\
Tương tác & Digital push button & Trạng thái nhấn/thả, thường là tín hiệu logic\\
Áp suất & Pressure sensor & Áp suất khí/quá trình, suy ra cao độ trong điều kiện phù hợp\\\bottomrule
\end{tabularx}
\end{table}

Ví dụ cảm biến trong slide gồm: nhiệt độ, ẩm đất, độ ẩm không khí, hồng ngoại, khí gas, CO$_2$, nút nhấn, alcohol, siêu âm, LDR, rain, photodiode, proximity, PIR, gyroscope. Cần phân biệt \emph{nguyên lý} và \emph{board/module}: cùng là cảm biến nhiệt độ nhưng có module analog, digital và cảm biến tiếp xúc/không tiếp xúc khác nhau.
\sourcepages{5--6, 11--12}

\section{Kết nối cảm biến với vi điều khiển (mở rộng kỹ thuật)}
\subsection{ADC, lượng tử hóa và ví dụ tính toán}
ADC $N$ bit tạo ra $2^N$ mức mã, thông thường từ $0$ đến $2^N-1$. Với ADC đơn cực lý tưởng tham chiếu $V_{\mathrm{ref}}$ và điện áp hợp lệ trong $[0,V_{\mathrm{ref}}]$, công thức quy đổi xấp xỉ từ mã $D$ là
\begin{equation}\boxed{V_{\mathrm{in}}\approx\frac{D}{2^N-1}V_{\mathrm{ref}}}\qquad (0\le D\le 2^N-1).\end{equation}
Bước lượng tử hóa danh định thường viết $\Delta V\approx V_{\mathrm{ref}}/2^N$; một số cách quy đổi cuối thang dùng $2^N-1$. Phải tuân theo quy ước của ADC và thư viện cụ thể. Ví dụ ADC 10 bit, $V_{\mathrm{ref}}=5\,\mathrm{V}$, nhận $D=512$: $V_{\mathrm{in}}\approx512\times5/1023\approx2.50\,\mathrm{V}$. Đây là \textbf{độ phân giải mã}, không nói trước sai số hệ thống.

\subsection{Cảm biến điện trở: mạch phân áp}
Một phần tử cảm biến thay đổi điện trở, chẳng hạn LDR hay NTC, cần được đo dưới dạng điện áp. Với $R_1$ nối từ nguồn $V_{CC}$ đến nút đo, $R_2$ nối từ nút đo xuống GND:
\begin{equation}V_{\mathrm{out}}=V_{CC}\frac{R_2}{R_1+R_2}.\end{equation}
\begin{figure}[htbp]\centering
\begin{tikzpicture}[font=\small\sffamily,>=Latex]
\node at (0,3.5){$V_{CC}$}; \draw[thick] (0,3.2)--(0,2.8);
\draw[draw=blueiot,thick,fill=paleblue](-.32,1.9) rectangle(.32,2.8);
\node[right] at (.4,2.35){$R_1$ cố định};
\draw[thick] (0,1.9)--(0,1.35)--(2.3,1.35);
\draw[->,tealiot,thick] (2.3,1.35)--(3.2,1.35) node[right]{$V_{out}$ đến ADC};
\draw[draw=blueiot,thick,fill=paleteal](-.32,.3) rectangle(.32,1.15);
\draw[thick] (0,1.35)--(0,1.15); \node[right] at (.4,.75){$R_2$ (LDR/NTC)};
\draw[thick] (0,.3)--(0,0); \draw[thick](-.5,0)--(.5,0); \draw[thick](-.3,-.13)--(.3,-.13); \draw[thick](-.14,-.26)--(.14,-.26);
\end{tikzpicture}
\caption{Đấu nối phân áp nguyên lý. Kiểm tra điện áp nguồn trước khi đưa vào ADC.}
\end{figure}

\subsection{I\textsuperscript{2}C, SPI, UART và ngõ vào số}
\begin{table}[htbp]\centering\small
\caption{Khi nào dùng giao tiếp nào?}
\begin{tabularx}{\textwidth}{p{18mm}p{42mm}X}
\toprule
Giao tiếp & Dây và nguyên lý & Lưu ý thiết kế\\\midrule
GPIO số & Một chân đọc HIGH/LOW, ngắt cạnh & Cần điện trở kéo lên/kéo xuống phù hợp; chống dội nút nhấn\\
ADC & Một chân analog tham chiếu mức điện áp & Chống aliasing, lọc nhiễu, hiệu chuẩn, không quá áp\\
I\textsuperscript{2}C & SDA/SCL, địa chỉ thiết bị trên bus & Điện trở kéo lên, địa chỉ trùng, điện áp logic, chiều dài bus\\
SPI & SCLK, MOSI, MISO, CS & Nhanh, nhiều dây; mỗi slave thường cần CS\\
UART & TX/RX, mức baud, định dạng bit & Nối chéo TX-RX, GND chung, không nhầm TTL với RS-232\\
1-Wire & Dây dữ liệu cùng GND, giao thức định thời & Điện trở kéo lên; thời gian và chiều dài đường truyền\\\bottomrule
\end{tabularx}
\end{table}

\section{Mô hình nút cảm biến IoT}
Theo sơ đồ trang 16--17 slide, nút IoT gồm \emph{Sensors + Board + Connectivity} và được phân rã thành cảm biến, mạch tiền xử lý/ADC, vi điều khiển, bộ nhớ, truyền thông có dây/không dây và khối nguồn. Lớp ứng dụng còn cần firmware, định danh và bảo mật.

\begin{figure}[htbp]\centering
\begin{tikzpicture}[font=\small\sffamily,node distance=7mm]
\node[block] (sens) {Cảm biến\\analog / digital};
\node[block,right=of sens] (front) {Analog front-end\\ADC / bus};
\node[block,right=of front] (cpu) {Vi điều khiển\\CPU + RAM + Flash};
\node[block,right=of cpu] (net) {Truyền thông\\Wi-Fi / BLE / RF};
\node[block,below=of cpu] (store) {Lưu trữ / nhật ký\\Flash / thẻ nhớ};
\node[block,below=of front] (power) {Nguồn pin / USB\\Sạc -- ổn áp -- bảo vệ};
\draw[arrow](sens)--(front); \draw[arrow](front)--(cpu); \draw[arrow](cpu)--(net);
\draw[<->,thick,blueiot] (cpu)--(store);
\draw[arrow] (power)--(front); \draw[arrow] (power.east) -- (store.west);
\end{tikzpicture}
\caption{Sơ đồ khối cấu trúc nút IoT tái dựng từ nội dung bài giảng.}
\end{figure}
\sourcepages{16--17}

\subsection{Nguồn và độ tin cậy đo}
Thiết kế nút cảm biến cần lưu ý dòng khởi động và dòng đỉnh của RF, mức nghỉ của cảm biến, chu kỳ đo, khả năng giảm công suất khi ngủ và tính toàn vẹn dữ liệu. Phương pháp ước lượng cơ bản:
\begin{equation}I_{\mathrm{avg}}=\sum_i D_i I_i,\qquad T_{\mathrm{h}}\approx\frac{C_{\mathrm{usable}}\,(\mathrm{mAh})}{I_{\mathrm{avg}}\,(\mathrm{mA})}.\end{equation}
Đây là ước tính lý tưởng theo dung lượng \emph{khả dụng}, chưa tính hiệu suất DC/DC, tự xả, suy giảm pin và nhiệt độ. Chu kỳ lấy mẫu nên dựa trên dải tần tín hiệu, yêu cầu trễ, tuổi thọ pin và chi phí truyền thông; lấy mẫu nhanh hơn không mặc nhiên làm dữ liệu chính xác hơn.

\section{Ví dụ thiết kế: đo nhiệt độ và điều khiển quạt}
Một cảm biến đọc nhiệt độ theo chu kỳ, MCU lọc trung bình trượt, so sánh với hai ngưỡng và xuất tín hiệu đến mạch công suất điều khiển quạt. Hai ngưỡng $T_{\mathrm{on}}>T_{\mathrm{off}}$ tạo \emph{hysteresis}, tránh đóng ngắt liên tục. Nếu sensor báo lỗi, hệ thống cần chính sách rõ ràng: tắt an toàn, báo lỗi hoặc bật quạt phòng ngừa tùy nguy cơ.
\begin{tipbox}
\textbf{Ví dụ giả định:} bật quạt khi $T\ge 32^\circ$C, tắt khi $T\le 29^\circ$C. Mức này chỉ minh họa thuật toán, không phải nhiệt độ khuyến nghị cho một loại thiết bị cụ thể. MCU không được trực tiếp cấp dòng cho mô-tơ.
\end{tipbox}

\section{Câu hỏi ôn tập và bài tập}
\begin{enumerate}
\item Trình bày $s=f(m)$, phân biệt tín hiệu đầu ra của cảm biến với đại lượng vật lý cần đo. Nêu một ví dụ có đường truyền xử lý hoàn chỉnh.
\item Theo phân loại của bài giảng, giải thích tại sao radar và sonar được coi là chủ động, cặp nhiệt điện là thụ động; nhận xét việc dùng từ ``không cần năng lượng''.
\item So sánh cảm biến analog với cảm biến digital; giải thích vì sao một cảm biến có giao tiếp I\textsuperscript{2}C có thể xuất nhiều giá trị đo dù tín hiệu trên dây là logic 0/1.
\item Tính mã ADC 10 bit gần đúng khi điện áp vào $1.65$ V và $V_{ref}=3.3$ V. \emph{Đáp số:} gần $512$.
\item Với phân áp $R_1=10\,\mathrm{k}\Omega$, $R_2=5\,\mathrm{k}\Omega$ và $V_{CC}=3.3$ V, tính $V_{out}$. \emph{Đáp số:} $1.1$ V.
\item Vẽ lại nút IoT dùng cảm biến nhiệt độ số, ESP32, Wi-Fi, nguồn pin và cảnh báo LED; ghi bus, điện áp logic, hướng truyền dữ liệu.
\item Phân tích ảnh hưởng của nhiệt độ, nhiễu nguồn, vị trí lắp đặt, sai số hiệu chuẩn đến độ tin cậy phép đo.
\end{enumerate}

\begin{notebox}\textbf{Tóm tắt cần nhớ:} Cảm biến biến đổi đại lượng thực thành tín hiệu; hiểu đúng đơn vị, dải đo, sai số và điện áp đầu ra quan trọng hơn việc ghi nhớ tên module. Ba trục phân loại (kích thích, hiện tượng, dạng đầu ra) có thể chồng lấp. Một nút IoT hoàn chỉnh cần cảm biến, mạch ghép nối, xử lý, nguồn và truyền thông.\end{notebox}
